\section{First Order Formulation}
Consider the following general first order formulation:
\begin{gather}
    \left\{\begin{array}{l}
       \ddot{u}+c_1u+ c_2f=0,\\
        \dot{f}+c_3f=c_4\ddot{u}+c_5\dot{u}+c_6u.
    \end{array}\right.
\end{gather}
Without complicating the discussion, we here assume constants $c_n\in\mathbb{R}$.
Rewriting the above in the state space, one has
\begin{gather}
    \begin{bmatrix}
        \dot{u}\\
        \ddot{u}\\
        \dot{f}
    \end{bmatrix}=\begin{bmatrix}
      0  &1&0\\
       -c_1 &0&-c_2\\
    c_6-c_4c_1   &c_5&-c_3-c_4c_2
    \end{bmatrix}\begin{bmatrix}
        u\\
        \dot{u}\\
        f
    \end{bmatrix}.
\end{gather}
The Routh--Hurwitz criterion can be expressed as
\begin{gather}
    c_2c_4+c_3>0,\quad
    c_1+c_2c_5>0,\quad
    c_1c_3+c_2c_6>0,\quad
    c_2\left(c_1c_4+c_2c_4c_5+c_3c_5-c_6\right)>0.
\end{gather}
\subsection{$c_5=c_6=0$}
The above simplifies to
\begin{gather}
    c_2c_4+c_3>0,\quad
    c_1>0,\quad
    c_1c_3>0,\quad
    c_1c_2c_4>0.
\end{gather}
Since $c_1>0$ and $c_1c_3>0$, one has $c_3>0$.
It is thus equivalent to
\begin{gather}
    \boxed{c_1>0,\quad
    c_3>0,\quad
    c_2c_4>0.}
\end{gather}
\subsection{$c_4=c_6=0$}
The above simplifies to
\begin{gather}
     c_3>0,\quad
    c_1+c_2c_5>0,\quad
    c_1c_3>0,\quad
    c_2c_3c_5>0.
\end{gather}
Since $c_1c_3>0$ and $c_3>0$, one has $c_1>0$.
It is thus equivalent to
\begin{gather}
    \boxed{c_1>0,\quad
    c_3>0,\quad
    c_2c_5>0.}
\end{gather}
\subsection{$c_4=c_5=0$}
The above simplifies to
\begin{gather}
    \boxed{
    c_3>0,\quad
    c_1>0,\quad
    c_1c_3+c_2c_6>0,\quad
    c_2c_6<0.}
\end{gather}